% !TeX root = ../../Book4.tex
% 纲要标题：### 7.4 特殊消解
% 纲要位置：计划纲要.md，第 2640 行。

\section{特殊消解}
\label{b4:sec:0704}


% 纲要内容（仅作写作计划，尚非教材正文）：
%
% * 主理想整环模的短消解
% * Koszul 消解简介
% * 群代数和代数的 bar 消解预告
%
% ---
%

\subsection{主理想整环模的短消解}
\label{b4:subsec:0704:01}

第二卷定理4.3.1证明了主理想整环上有限秩自由模的子模自由。为了处理并非有限生成的模，需要把同一思路延伸到任意基集合。

\begin{theorem}\label{b4:thm:pid-free-submodule}
设 \(R\) 是主理想整环，\(F\) 是任意秩自由 \(R\) 模。每个子模 \(N\subseteq F\) 都自由。因此每个 \(R\) 模 \(M\) 都有长度至多为一的自由消解 \(0\to F_1\to F_0\to M\to0\)。
\end{theorem}
\begin{proof}
用选择公理良序排列 \(F\) 的一组基，写成 \((e_\alpha)_{\alpha<\kappa}\)。令 \(F_\beta\) 由 \(\alpha<\beta\) 的基向量生成，\(N_\beta=N\cap F_\beta\)。坐标投影使 \(N_{\alpha+1}/N_\alpha\) 同构于 \(R\) 的一个理想。若它非零，取其生成元 \(a_\alpha\)，并在 \(N_{\alpha+1}\) 选一个 \(e_\alpha\) 坐标为 \(a_\alpha\) 的元素 \(x_\alpha\)。任意 \(y\in N_{\alpha+1}\) 的该坐标是 \(ra_\alpha\)，所以 \(y-rx_\alpha\in N_\alpha\)；而 \(rx_\alpha\in N_\alpha\) 蕴含 \(ra_\alpha=0\)，由整环性得 \(r=0\)。因此 \(N_{\alpha+1}=N_\alpha\oplus Rx_\alpha\)。若该理想为零，则 \(N_{\alpha+1}=N_\alpha\)。

极限序数 \(\lambda\) 处有 \(N_\lambda=\bigcup_{\beta<\lambda}N_\beta\)，因为每个向量只含有限个非零坐标。超限归纳遂表明全部选出的 \(x_\alpha\) 构成 \(N\) 的基：生成性由上述分解给出，任一有限关系则由最大下标的非零坐标逐次消去。最后取自由满射 \(F_0\to M\)，其核自由，即得短消解。
\end{proof}

有限生成 \(M\) 时可先取有限秩 \(F_0\)，并由有限秩版本得到有限秩 \(F_1\)。若 \(M=R/(a)\)、\(a\ne0\)，其短消解就是乘 \(a\) 的映射；整环假设在这里保证该映射单射。

\subsection{Koszul 消解简介}
\label{b4:subsec:0704:02}

对一个元素 \(x\in R\)，两项复形 \(K(x)=(0\to R\xrightarrow{x}R\to0)\) 可以是消解，也可以不是。其零次同调为 \(R/(x)\)，一次同调为 \(\{\,r\in R\mid xr=0\,\}\)。因此正合性准确地要求 \(x\) 不是零因子。

\begin{definition}\label{b4:def:koszul-preview}
设 \(R\) 是交换环，\(x_1,\ldots,x_r\in R\)。取自由模 \(R^r\) 的基 \(e_1,\ldots,e_r\)，令 \(K_n=\bigwedge^nR^r\)，并定义
\[
d(e_{i_1}\wedge\cdots\wedge e_{i_n})
=\sum_{j=1}^n(-1)^{j-1}x_{i_j}
 e_{i_1}\wedge\cdots\widehat{e_{i_j}}\cdots\wedge e_{i_n}.
\]
所得链复形记为 \(K_\bullet(x_1,\ldots,x_r;R)\)，称为\term[koszul-fuxing]{Koszul 复形}[Koszul complex]。
\end{definition}

删去两个基向量的两种次序符号相反，且 \(x_ix_j=x_jx_i\)，所以 \(d^2=0\)。外幂及楔积的定义与基已在第二卷第九章建立。第三卷定理15.3.2已把正则列的 Koszul 消解作为交换代数中的主要结果处理；下面给出基于本卷同调长正合列的第二证明，只证明消解所需的方向，不在这里重做第三卷的深度理论。

\begin{proposition}\label{b4:prop:koszul-resolution}
设 \(R\) 是交换环，\(x_1,\ldots,x_r\) 满足 \((x_1,\ldots,x_r)\ne R\)，且对每个 \(i\)，乘 \(x_i\) 在 \(R/(x_1,\ldots,x_{i-1})\) 上单射。则 \(K_\bullet(x_1,\ldots,x_r;R)\) 是 \(R/(x_1,\ldots,x_r)\) 的有限自由消解。
\end{proposition}
\begin{proof}
\(r=0\) 时取集中于零次的 \(R\)，结论成立。对正整数 \(r\) 归纳。\(r=1\) 就是前面的两项计算。设 \(K'=K_\bullet(x_1,\ldots,x_{r-1};R)\)。按含不含 \(e_r\) 分解楔积，每次 \(K_n=K'_n\oplus(K'_{n-1}\wedge e_r)\)，且
\[
d(a+b\wedge e_r)=d'a+(-1)^{n-1}x_rb+(d'b)\wedge e_r.
\]
故有短正合复形列，左端为 \(K'\)，商复形在次数 \(n\) 为 \(K'_{n-1}\)，微分为 \(d'\)。其连接态射在同调上是乘 \(x_r\)，差至多为显式的符号 \((-1)^{n-1}\)。归纳假设使 \(K'\) 只在零次有同调。因此长正合列给出 \(H_n(K)=0\) 对 \(n\ge2\)，并给出
\[
0\longrightarrow H_1(K)\longrightarrow R/(x_1,\ldots,x_{r-1})
\xrightarrow{x_r}R/(x_1,\ldots,x_{r-1})
\longrightarrow H_0(K)\longrightarrow0.
\]
假设使中间映射单射，故一次同调也消失，零次同调为所需商环。
\end{proof}

例如 \(k[x,y]\) 上的 \(k=R/(x,y)\) 有消解
\[
0\longrightarrow R\xrightarrow{c\mapsto(-yc,xc)}R^2
\xrightarrow{(a,b)\mapsto xa+yb}R\longrightarrow k\longrightarrow0.
\]
它在下一章计算 Tor 时会产生一次和二次的不消失项。

\subsection{群代数与代数的 bar 消解}
\label{b4:subsec:0704:03}

群环和一般代数的自然模往往没有短消解，却有一种由“插入一个符号”控制正合性的自由消解。这里介绍其形式；后面群上同调与 Hochschild 上同调的章节将说明它为何适合相应问题。

\begin{example}\label{b4:exm:homogeneous-bar}
设 \(G\) 是群，\(k\) 是交换环。令 \(B_n\) 为以 \((g_0,\ldots,g_n)\in G^{n+1}\) 为基的自由 \(k\) 模，并让 \(G\) 同时左乘所有坐标。它是自由左 \(kG\) 模：每个轨道唯一含有首坐标为 \(1\) 的元组。令
\[
d(g_0,\ldots,g_n)=\sum_{i=0}^n(-1)^i(g_0,\ldots,\widehat{g_i},\ldots,g_n),
\qquad\varepsilon(g_0)=1.
\]
两次删除的项成对抵消，故 \(d^2=0\)。在底层 \(k\) 模上令 \(s_n(g_0,\ldots,g_n)=(1,g_0,\ldots,g_n)\)、\(s_{-1}(1)=(1)\)，则 \(ds+sd=1\)，因为删除新插入的首项给出原元组，其余项与 \(sd\) 抵消。因此增广复形正合，是平凡左 \(kG\) 模 \(k\) 的自由消解。\(s\) 通常不具 \(G\) 等变性；这不妨碍正合，却不证明它作为 \(kG\) 复形可缩。
\end{example}

代数的 bar 构造把元组换成张量串。例如域 \(k\) 上含幺代数 \(A\) 的各项 \(A^{\otimes_k(n+2)}\)，以首尾因子承担 \(A\) 双模作用，微分交错地乘合相邻因子，增广为 \(A\otimes_k A\to A\) 的乘法。在左端插入 \(1\) 给出底层向量空间的同样收缩。每项作为 \(A\otimes_k A^{\mathrm{op}}\) 模自由，因为中间张量因子有 \(k\) 基。此处只预告这一形式及其正合原因，后续章节将详细建立双模约定、归一化和所导出的上同调计算。
